% ============================================================
%  Bayesian UQ for Porous Media Permeability Inversion
%  Full derivations + figures
% ============================================================
\documentclass[11pt, a4paper]{article}

\usepackage{amsmath, amssymb, amsthm}
\usepackage{graphicx}
\usepackage[margin=2.5cm]{geometry}
\usepackage{hyperref}
\usepackage{booktabs}
\usepackage{xcolor}
\usepackage{subcaption}
\usepackage{enumitem}
\usepackage{mathtools}
\usepackage{float}
\usepackage[ruled, linesnumbered]{algorithm2e}
\usepackage{parskip}
\usepackage{tcolorbox}

\hypersetup{colorlinks=true, linkcolor=blue!60!black, citecolor=green!50!black}

% ─── notation shortcuts ───────────────────────────────────────────────
\newcommand{\R}{\mathbb{R}}
\newcommand{\N}{\mathcal{N}}
\newcommand{\E}{\mathbb{E}}
\newcommand{\Var}{\operatorname{Var}}
\newcommand{\IG}{\operatorname{InvGamma}}
\newcommand{\Ga}{\operatorname{Gamma}}
\newcommand{\GP}{\mathcal{GP}}
\newcommand{\norm}[1]{\left\|#1\right\|}
\newcommand{\abs}[1]{\left|#1\right|}
\newcommand{\bxi}{\boldsymbol{\xi}}   % bold KL coefficient vector

% ─── theorem-like environments ────────────────────────────────────────
\theoremstyle{definition}
\newtheorem{definition}{Definition}[section]
\newtheorem{proposition}{Proposition}[section]
\newtheorem{remark}{Remark}[section]

% ─── figure path ─────────────────────────────────────────────────────
\graphicspath{{./figures/}{./conjugate_priors/figures/}}

% =====================================================================
\begin{document}

\title{\textbf{Bayesian Uncertainty Quantification for Permeability Inversion
in Porous Media}\\[6pt]
\large Probability Distributions, Conjugate Priors, Gaussian Random Fields,
and pCN-Gibbs Sampling --- Full Derivations with Figures}
\author{}
\date{\today}
\maketitle
\tableofcontents
\newpage

% =====================================================================
\section{Introduction}
% =====================================================================

Estimating the permeability field $k(x)$ of a porous medium from sparse
pressure measurements is a canonical inverse problem in subsurface
engineering.  The forward map is Darcy's law, which is well-posed: given
$k(x)$, compute $p(x)$.  The inverse direction---recovering $k(x)$ from
a handful of noisy pressure readings---is \emph{ill-posed}: many permeability
fields can reproduce the same observations within measurement error.

The Bayesian framework converts this into a well-posed problem by placing a
probability distribution over $k(x)$ and using data to update it.  The result
is not a single ``best'' permeability map but a \emph{posterior distribution}
that quantifies how uncertain we remain after the experiment.

This document derives every piece of machinery used in the accompanying Python
code, starting from the probability distributions introduced in UQ textbooks
(Normal, Gamma, Inverse-Gamma) and building up to the full pCN-Gibbs posterior
sampler.

% =====================================================================
\section{The Forward Problem: 1D Single-Phase Darcy Flow}
% =====================================================================

\subsection{Governing equation}

We consider single-phase, steady, incompressible flow on $[0,1]$ governed by
Darcy's law with no volumetric source:
\begin{equation}
  -\frac{\mathrm{d}}{\mathrm{d}x}\!\left[k(x)\,\frac{\mathrm{d}p}{\mathrm{d}x}\right] = 0,
  \qquad x\in(0,1),
  \label{eq:darcy}
\end{equation}
with Dirichlet boundary conditions $p(0)=p_L$ and $p(1)=p_R$.  The
permeability $k(x)>0$ is the \emph{unknown field} we wish to infer.

\subsection{Analytic solution}

Integrating \eqref{eq:darcy} once gives $k(x)\,p'(x)=C$ (constant flux).
A second integration yields:
\begin{equation}
  \boxed{%
  p(x) = p_L + (p_R - p_L)\,\frac{R(x)}{R(1)},
  \quad
  R(x) \coloneqq \int_0^x \frac{\mathrm{d}s}{k(s)}}
  \label{eq:darcy_sol}
\end{equation}
where $R(x)$ is the cumulative hydraulic resistance.  The Darcy flux
(constant throughout the domain) is
\begin{equation}
  q = \frac{p_L - p_R}{R(1)}.
\end{equation}
Equation \eqref{eq:darcy_sol} is computed numerically via the
cumulative trapezoidal rule on the $N$-point grid
$\{x_j\}_{j=0}^{N-1}$, giving an exact (to quadrature order) evaluation
in $O(N)$ operations.

\subsection{Observation operator}

Let $\mathbf{x}_s = (x_{s_1},\ldots,x_{s_m})$ be the sensor locations.
The observation operator $H$ extracts the pressure at these points:
\begin{equation}
  H : k(\cdot) \mapsto \mathbf{d} = \bigl[p(x_{s_1}),\ldots,p(x_{s_m})\bigr]^\top
  \in \R^m.
\end{equation}
Composing with the forward solve we write $G=H\circ\text{Darcy}$, so that
$\mathbf{d}_{\mathrm{pred}} = G(k)$.

% =====================================================================
\section{Probability Distributions in Bayesian UQ}
% =====================================================================

Every UQ textbook opens with the Normal, Gamma, and Inverse-Gamma
distributions.  Their relevance is not axiomatic---each one earns its place
by solving a specific subproblem.

\begin{figure}[H]
  \centering
  \includegraphics[width=\textwidth]{fig1_distributions}
  \caption{The three foundational distributions for Gaussian-likelihood UQ.
    \textit{Left}: Normal---models measurement noise and the log-permeability
    prior.
    \textit{Centre}: Inverse-Gamma---conjugate hyperprior on the noise
    variance $\sigma_\varepsilon^2$.
    \textit{Right}: Gamma---conjugate prior on the precision
    $\tau=1/\sigma_\varepsilon^2$.}
  \label{fig:distributions}
\end{figure}

\subsection{Normal (Gaussian) distribution}

\begin{definition}
  $X\sim\N(\mu,\sigma^2)$ has pdf
  \[
    p(x;\mu,\sigma^2) = \frac{1}{\sqrt{2\pi\sigma^2}}
    \exp\!\left(-\frac{(x-\mu)^2}{2\sigma^2}\right).
  \]
\end{definition}

The Normal distribution serves three distinct roles in our problem:
\begin{enumerate}
  \item \textbf{Likelihood / measurement noise.}
    Observations are modelled as
    $\mathbf{d}_{\mathrm{obs}} = G(k) + \boldsymbol{\varepsilon}$,
    $\boldsymbol{\varepsilon}\sim\N(\mathbf{0},\sigma_\varepsilon^2 I)$.
    This is the maximum-entropy distribution consistent with a known
    variance, and it converts the data misfit $\norm{\mathbf{d}_{\mathrm{obs}}-G(k)}^2$
    into a log-likelihood.

  \item \textbf{Prior on the log-permeability field.}
    Because $k(x)>0$, we model $Y(x)=\log k(x)$ as a Gaussian random
    field $Y\sim\GP(\mu_Y, C)$, which guarantees $k=\exp(Y)>0$ by
    construction.

  \item \textbf{KL coefficients.}
    The Karhunen-Lo\`eve expansion (Section~\ref{sec:kl}) parametrises
    $Y(x)$ through coefficients $\xi_i\sim\N(0,1)$, reducing the
    infinite-dimensional inference to a problem in $\R^M$.
\end{enumerate}

\subsection{Gamma distribution}

\begin{definition}
  $X\sim\Ga(\alpha,\beta)$ with rate $\beta$ has pdf
  \[
    p(x;\alpha,\beta) = \frac{\beta^\alpha}{\Gamma(\alpha)}\,
    x^{\alpha-1}\,e^{-\beta x}, \quad x>0,
  \]
  with $\E[X]=\alpha/\beta$, $\Var[X]=\alpha/\beta^2$.
\end{definition}

The Gamma is the \emph{conjugate prior} for the precision
$\tau=1/\sigma_\varepsilon^2$ of a Gaussian likelihood: if
$\tau\sim\Ga(\alpha_0,\beta_0)$ then, after $n$ observations with
sum of squared residuals $\mathrm{SS}$, the posterior is
$\tau|\mathbf{d}\sim\Ga(\alpha_0+n/2,\;\beta_0+\mathrm{SS}/2)$.

\subsection{Inverse-Gamma distribution}

\begin{definition}
  $X\sim\IG(\alpha,\beta)$ with scale $\beta$ has pdf
  \[
    p(x;\alpha,\beta) = \frac{\beta^\alpha}{\Gamma(\alpha)}\,
    x^{-\alpha-1}\,e^{-\beta/x}, \quad x>0,
  \]
  with $\E[X]=\beta/(\alpha-1)$ for $\alpha>1$ and
  $\mathrm{Mode}[X]=\beta/(\alpha+1)$.
\end{definition}

If $X\sim\IG(\alpha,\beta)$ then $1/X\sim\Ga(\alpha,\beta)$, so Gamma and
Inverse-Gamma are in a dual relationship.  We use $\IG$ as the conjugate
prior for $\sigma_\varepsilon^2$ because it is the natural distribution for
a variance parameter (positive, right-skewed, supported on $(0,\infty)$).

The key closed-form update is derived in Section~\ref{sec:nig}:
\begin{equation}
  \sigma_\varepsilon^2\sim\IG(\alpha_0,\beta_0)
  \;\;\Longrightarrow\;\;
  \sigma_\varepsilon^2\,|\,\mathbf{r}\sim
  \IG\!\left(\alpha_0+\tfrac{n}{2},\;\beta_0+\tfrac{\norm{\mathbf{r}}^2}{2}\right)
  \label{eq:ig_update}
\end{equation}
where $\mathbf{r}=\mathbf{d}_{\mathrm{obs}}-G(\xi)$ are the residuals.

\subsection{Log-Normal distribution and the permeability prior}

Because $k=e^Y$ and $Y\sim\N(\mu_Y,\sigma_Y^2)$, the permeability itself
follows a log-normal:
\[
  p(k;\mu_Y,\sigma_Y)
  = \frac{1}{k\,\sigma_Y\sqrt{2\pi}}
    \exp\!\left(-\frac{(\log k - \mu_Y)^2}{2\sigma_Y^2}\right),
  \quad k>0.
\]
The median permeability is $e^{\mu_Y}$ and the coefficient of variation
grows with $\sigma_Y$.  Working in log-space (i.e.\ modelling $Y=\log k$
with a Gaussian field) is both statistically natural and computationally
convenient: positivity is guaranteed by construction and the posterior
inherits the Gaussian structure of the prior.

% =====================================================================
\section{Gaussian Random Field Prior and KL Expansion}
\label{sec:kl}
% =====================================================================

\subsection{Gaussian random fields}

A Gaussian random field on $[0,1]$ is a stochastic process
$\{Y(x):x\in[0,1]\}$ such that every finite-dimensional marginal is
multivariate Gaussian.  It is fully specified by its mean function
$\mu_Y(x)=\E[Y(x)]$ and covariance function (kernel)
$C(x,x')=\mathrm{Cov}[Y(x),Y(x')]$.

We use the exponential (Ornstein-Uhlenbeck) kernel:
\begin{equation}
  C(x,x') = \sigma_Y^2\,\exp\!\left(-\frac{\abs{x-x'}}{\ell}\right),
  \label{eq:exp_kernel}
\end{equation}
where $\sigma_Y^2$ is the marginal variance and $\ell$ is the
\emph{correlation length}.  Small $\ell$ produces rough, rapidly-varying
fields; large $\ell$ produces smooth, slowly-varying ones.

\subsection{Karhunen-Lo\`eve expansion}

Let $\{x_j\}_{j=0}^{N-1}$ be the $N$-point discretisation grid.  Define
the $N\times N$ covariance matrix $\mathbf{C}$ with
$C_{ij}=C(x_i,x_j)$.  Its symmetric eigendecomposition gives
\[
  \mathbf{C} = \boldsymbol{\Phi}\,\mathrm{diag}(\boldsymbol{\lambda})\,\boldsymbol{\Phi}^\top,
  \quad \lambda_1\geq\lambda_2\geq\cdots\geq\lambda_N\geq 0,
\]
where $\boldsymbol{\phi}_i$ is the $i$-th eigenvector (KL mode).

Retaining only the $M$ leading eigenpairs gives the \emph{truncated KL expansion}:
\begin{equation}
  \boxed{%
  Y(x) \approx \mu_Y + \sum_{i=1}^{M}\sqrt{\lambda_i}\,\phi_i(x)\,\xi_i,
  \qquad \xi_i\overset{\mathrm{i.i.d.}}{\sim}\N(0,1)}
  \label{eq:kl}
\end{equation}

Equation \eqref{eq:kl} converts the infinite-dimensional field inference
into a finite-dimensional problem: \emph{infer the $M$ scalar coefficients
$\xi=(\xi_1,\ldots,\xi_M)^\top$}.  The fraction of variance captured is
\[
  f_M = \frac{\sum_{i=1}^M\lambda_i}{\sum_{i=1}^N\lambda_i}.
\]
For the exponential kernel with $\ell=0.3$ and $N=65$, the first $M=20$
modes capture approximately 96.8\% of the total field variance.

\begin{figure}[H]
  \centering
  \includegraphics[width=\textwidth]{fig2_kl_prior}
  \caption{Karhunen-Lo\`eve expansion of the log-permeability GRF.
    \textit{Left}: eigenvalue decay of the exponential kernel;
    the rapid fall justifies truncation at $M=20$.
    \textit{Centre}: first five eigenfunctions (spatial fingerprints
    of each $\xi_i$).
    \textit{Right}: 15 prior sample fields drawn by sampling
    $\xi\sim\N(\mathbf{0},I)$ and applying \eqref{eq:kl}.}
  \label{fig:kl}
\end{figure}

\subsection{Wiener-Askey connection to polynomial chaos}

When a surrogate-based approach (polynomial chaos expansion, PCE) is
preferred over MCMC, the distribution chosen for the KL coefficients
dictates the orthogonal polynomial basis:
\begin{center}
\begin{tabular}{lll}
\toprule
Input distribution & Polynomial basis & Support \\
\midrule
$\xi_i\sim\N(0,1)$       & Hermite   & $(-\infty,\infty)$ \\
$\xi_i\sim\Ga(\alpha,1)$ & Laguerre  & $(0,\infty)$ \\
$\xi_i\sim\mathrm{Unif}(-1,1)$ & Legendre & $[-1,1]$ \\
$\xi_i\sim\mathrm{Beta}(\alpha,\beta)$ & Jacobi & $[0,1]$ \\
\bottomrule
\end{tabular}
\end{center}
Since we use $\xi_i\sim\N(0,1)$, the natural PCE basis is Hermite
polynomials.

% =====================================================================
\section{Bayesian Formulation of the Inverse Problem}
% =====================================================================

\subsection{Bayes' theorem}

Let $\xi\in\R^M$ be the KL coefficients (our inferred quantity) and
$\sigma_\varepsilon^2>0$ the noise variance (treated as unknown,
hierarchical model).  The observation model is
\begin{equation}
  \mathbf{d}_{\mathrm{obs}} = G(\xi) + \boldsymbol{\varepsilon},
  \qquad
  \boldsymbol{\varepsilon}\sim\N(\mathbf{0},\sigma_\varepsilon^2 I_m),
  \label{eq:obs_model}
\end{equation}
where $G(\xi)$ maps KL coefficients to predicted sensor pressures
(via the KL expansion and the Darcy solve).

Bayes' theorem gives the posterior:
\begin{equation}
  \underbrace{p(\xi,\sigma_\varepsilon^2\,|\,\mathbf{d}_{\mathrm{obs}})}_{\text{posterior}}
  \;\propto\;
  \underbrace{p(\mathbf{d}_{\mathrm{obs}}\,|\,\xi,\sigma_\varepsilon^2)}_{\text{likelihood}}
  \;\cdot\;
  \underbrace{p(\xi)}_{\text{field prior}}
  \;\cdot\;
  \underbrace{p(\sigma_\varepsilon^2)}_{\text{noise hyperprior}}
  \label{eq:bayes}
\end{equation}

\begin{figure}[H]
  \centering
  \includegraphics[width=\textwidth]{fig3_problem_setup}
  \caption{Synthetic experiment setup.
    \textit{Left}: true log-permeability field (one KL realisation)
    overlaid on the prior $\pm2\sigma_Y$ envelope.
    \textit{Right}: true pressure field and the 7 noisy sensor
    observations that constitute $\mathbf{d}_{\mathrm{obs}}$.}
  \label{fig:setup}
\end{figure}

\subsection{Three terms in the posterior}

\paragraph{Gaussian likelihood.}
\begin{equation}
  \log p(\mathbf{d}_{\mathrm{obs}}\,|\,\xi,\sigma_\varepsilon^2)
  = -\frac{m}{2}\log(2\pi\sigma_\varepsilon^2)
    -\frac{\norm{\mathbf{d}_{\mathrm{obs}}-G(\xi)}^2}{2\sigma_\varepsilon^2}
  \label{eq:loglik}
\end{equation}
This is the ``comparison to experiment'' term: it penalises the squared
misfit between the forward model prediction and the observed pressures,
weighted by the noise level.

\paragraph{KL-coefficient prior.}
The KL parametrisation absorbs the GRF covariance into the eigenvectors,
so the prior on coefficients is simply:
\begin{equation}
  \xi\sim\N(\mathbf{0},I_M)
  \;\Longrightarrow\;
  \log p(\xi) = -\tfrac{1}{2}\norm{\xi}^2 + \mathrm{const}
  \label{eq:xi_prior}
\end{equation}

\paragraph{Inverse-Gamma hyperprior.}
The noise variance $\sigma_\varepsilon^2$ is unknown.  We place a
weakly-informative conjugate prior:
\begin{equation}
  \sigma_\varepsilon^2\sim\IG(\alpha_0,\beta_0),
  \quad
  \log p(\sigma_\varepsilon^2)
  = -(\alpha_0+1)\log\sigma_\varepsilon^2 - \frac{\beta_0}{\sigma_\varepsilon^2}
    + \mathrm{const}
  \label{eq:ig_prior}
\end{equation}

\subsection{Regularisation interpretation}

Dropping the noise hyperprior and fixing $\sigma_\varepsilon^2$, the
log-posterior is
\[
  \log p(\xi\,|\,\mathbf{d}_{\mathrm{obs}}) \propto
  -\frac{\norm{\mathbf{d}_{\mathrm{obs}}-G(\xi)}^2}{2\sigma_\varepsilon^2}
  -\frac{\norm{\xi}^2}{2}.
\]
Maximising this is exactly \emph{Tikhonov regularisation} with
regularisation parameter $\lambda=\sigma_\varepsilon^2$.  The Bayesian
framework additionally provides:
(i) a principled choice of $\lambda$ via the hyperprior;
(ii) credible intervals from the full posterior, not just the MAP point.

% =====================================================================
\section{Conjugate Priors: Theory and Full Derivations}
\label{sec:conjugate}
% =====================================================================

\begin{definition}
  A prior $p(\theta\,|\,\mathbf{h})$ is \textbf{conjugate} to the likelihood
  $p(\mathbf{x}\,|\,\theta)$ if the posterior $p(\theta\,|\,\mathbf{x})$
  belongs to the same parametric family as the prior, with updated
  hyperparameters $\mathbf{h}\to\mathbf{h}_n$.
\end{definition}

The consequence is that Bayesian updating reduces to a few algebraic
formulas---no numerical integration, no MCMC, no approximation.

\begin{figure}[H]
  \centering
  \includegraphics[width=\textwidth]{fig1_triptych}
  \caption{The prior-likelihood-posterior triptych for the Normal-Normal
    conjugate pair (Case 1).
    The prior (blue), likelihood function (orange), and posterior (green)
    are all Gaussians; only the mean and width change.}
  \label{fig:triptych}
\end{figure}

% ─── Case 1 ────────────────────────────────────────────────────────────
\subsection{Case 1: Normal-Normal (inferring the mean, $\sigma$ known)}
\label{sec:nn}

\paragraph{Model.}
\[
  x_i\,|\,\mu \overset{\mathrm{i.i.d.}}{\sim} \N(\mu,\sigma^2),\quad i=1,\ldots,n,
  \qquad
  \mu\sim\N(\mu_0,\tau_0^2).
\]

\paragraph{Derivation.}
The log-posterior is
\begin{align}
  \log p(\mu\,|\,\mathbf{x}) &\propto
  \sum_{i=1}^n \log\N(x_i;\mu,\sigma^2) + \log\N(\mu;\mu_0,\tau_0^2)
  \notag\\
  &= -\frac{1}{2}\left[
    \frac{\sum(x_i-\mu)^2}{\sigma^2}
    +\frac{(\mu-\mu_0)^2}{\tau_0^2}
  \right] + \mathrm{const}.
  \label{eq:nn_logpost}
\end{align}
Expanding $\sum(x_i-\mu)^2 = n\mu^2 - 2n\bar{x}\mu + \sum x_i^2$ and
collecting terms in $\mu$:
\begin{equation}
  \log p(\mu\,|\,\mathbf{x}) \propto
  -\frac{1}{2}\left[\underbrace{\left(\frac{n}{\sigma^2}+\frac{1}{\tau_0^2}\right)}_{\displaystyle 1/\tau_n^2}
  \mu^2
  - 2\underbrace{\left(\frac{n\bar{x}}{\sigma^2}+\frac{\mu_0}{\tau_0^2}\right)}_{\displaystyle \mu_n/\tau_n^2}\mu\right].
\end{equation}
Completing the square gives $\log p(\mu\,|\,\mathbf{x})\propto -(\mu-\mu_n)^2/(2\tau_n^2)$, so:

\begin{tcolorbox}[colback=blue!5, colframe=blue!50!black, title=Normal-Normal update, fontupper=\small]
Posterior: $\;\mu\,|\,\mathbf{x}\sim\N(\mu_n,\tau_n^2)$ \quad with
\begin{align}
  \frac{1}{\tau_n^2} &= \frac{1}{\tau_0^2} + \frac{n}{\sigma^2}
    \qquad\text{(precisions add)}\label{eq:nn_prec}\\[4pt]
  \mu_n &= \tau_n^2\!\left(\frac{\mu_0}{\tau_0^2}+\frac{n\bar{x}}{\sigma^2}\right)
    = \frac{(1/\tau_0^2)\,\mu_0 + (n/\sigma^2)\,\bar{x}}{1/\tau_0^2+n/\sigma^2}
    \quad\text{(precision-weighted mean)}\label{eq:nn_mean}
\end{align}
\end{tcolorbox}

\paragraph{Key insight.}
Equation~\eqref{eq:nn_prec} shows that the posterior \emph{precision}
$1/\tau_n^2$ grows \emph{linearly} in $n$; equivalently the variance
$\tau_n^2\sim 1/n$ for large $n$.  Adding data is like adding parallel
resistors: more resistors $\Rightarrow$ lower total resistance.

\paragraph{Pseudo-data interpretation.}
The prior $\N(\mu_0,\tau_0^2)$ can be viewed as coming from $n_0=\sigma^2/\tau_0^2$
imaginary observations all equal to $\mu_0$.  Real data updates the same
counter: $n_{\mathrm{eff}}=n_0+n$, $\bar{x}_{\mathrm{eff}}=(n_0\mu_0+n\bar{x})/n_{\mathrm{eff}}$.

\begin{figure}[H]
  \centering
  \includegraphics[width=\textwidth]{fig2_sequential_nn}
  \caption{Sequential Normal-Normal updating at $n=0,1,3,8,20,50$.
    The posterior sharpens around $\mu_{\mathrm{true}}=1.5$ (red dashed line)
    as data accumulates.  At every step the posterior is exactly Gaussian.}
  \label{fig:seq_nn}
\end{figure}

% ─── Case 2 ────────────────────────────────────────────────────────────
\subsection{Case 2: Normal-Inverse-Gamma (inferring $\sigma^2$, $\mu$ known)}
\label{sec:nig}

This is the case used directly in the Darcy Gibbs step.

\paragraph{Model.}
\[
  x_i\,|\,\sigma^2 \overset{\mathrm{i.i.d.}}{\sim}\N(\mu,\sigma^2),
  \quad
  \sigma^2\sim\IG(\alpha_0,\beta_0).
\]

\paragraph{Derivation.}
The InvGamma prior pdf is $p(\sigma^2)\propto (\sigma^2)^{-\alpha_0-1}e^{-\beta_0/\sigma^2}$.
The likelihood is
\[
  p(\mathbf{x}\,|\,\sigma^2)
  = (2\pi\sigma^2)^{-n/2}\exp\!\left(-\frac{\mathrm{SS}}{2\sigma^2}\right),
  \quad \mathrm{SS}=\sum_{i=1}^n(x_i-\mu)^2.
\]
The posterior is therefore
\begin{align}
  p(\sigma^2\,|\,\mathbf{x})
  &\propto (\sigma^2)^{-n/2}\exp\!\left(-\frac{\mathrm{SS}}{2\sigma^2}\right)
           \cdot(\sigma^2)^{-\alpha_0-1}e^{-\beta_0/\sigma^2} \notag\\
  &= (\sigma^2)^{-(\alpha_0+n/2)-1}
     \exp\!\left(-\frac{\beta_0+\mathrm{SS}/2}{\sigma^2}\right).
  \label{eq:nig_post}
\end{align}
Recognising \eqref{eq:nig_post} as an InvGamma kernel:

\begin{tcolorbox}[colback=orange!5, colframe=orange!60!black, title=Normal-InvGamma update, fontupper=\small]
Posterior: $\;\sigma^2\,|\,\mathbf{x}\sim\IG(\alpha_n,\beta_n)$ \quad with
\begin{align}
  \alpha_n &= \alpha_0 + \frac{n}{2} \label{eq:nig_alpha}\\
  \beta_n  &= \beta_0 + \frac{\mathrm{SS}}{2}
    = \beta_0 + \frac{1}{2}\sum_{i=1}^n(x_i-\mu)^2 \label{eq:nig_beta}
\end{align}
Posterior mean: $\;\E[\sigma^2\,|\,\mathbf{x}]=\beta_n/(\alpha_n-1)$
\quad (for $\alpha_n>1$).
\end{tcolorbox}

\paragraph{Information accumulation.}
$\alpha_n$ grows at rate $\frac{1}{2}$ per observation (each datum
contributes ``half a degree of freedom''), while $\beta_n$ accumulates
the sum of squared residuals.  Both updates are additive and can be
applied one datum at a time with identical results.

\paragraph{Application to the Darcy Gibbs step.}
With KL coefficients $\xi$ fixed, the Darcy residuals
$\mathbf{r}=\mathbf{d}_{\mathrm{obs}}-G(\xi)$ play the role of $\{x_i-\mu\}$,
and $\mathrm{SS}=\norm{\mathbf{r}}^2$.  Equation~\eqref{eq:ig_update}
gives a \emph{single closed-form sample} for $\sigma_\varepsilon^2$
without any accept/reject step.

\begin{figure}[H]
  \centering
  \includegraphics[width=\textwidth]{fig3_sequential_nig}
  \caption{Sequential Normal-InvGamma updating at $n=0,1,3,8,20,50$.
    The InvGamma posterior concentrates near $\sigma^2_{\mathrm{true}}=0.16$
    (red dashed).  Throughout, the posterior remains an InvGamma distribution.}
  \label{fig:seq_nig}
\end{figure}

% ─── Precision geometry ───────────────────────────────────────────────

\begin{figure}[H]
  \centering
  \includegraphics[width=\textwidth]{fig4_precision}
  \caption{Information geometry of conjugate updates.
    \textit{Left}: posterior variance $\tau_n^2$ decreases hyperbolically in $n$.
    \textit{Centre}: posterior precision $1/\tau_n^2$ is linear in $n$
    (information adds).
    \textit{Right}: InvGamma shape parameter $\alpha_n=\alpha_0+n/2$ is
    also linear (``half-observation'' accumulation).}
  \label{fig:precision}
\end{figure}

% ─── Case 3 ────────────────────────────────────────────────────────────
\subsection{Case 3: Normal-Normal-InvGamma (joint inference of $\mu$ and $\sigma^2$)}
\label{sec:nig_joint}

\paragraph{Model.}
The joint Normal-InvGamma (NIG) prior factorises as
\begin{equation}
  p(\mu,\sigma^2) = \N\!\left(\mu\,\Big|\,\mu_0,\frac{\sigma^2}{\kappa_0}\right)
  \cdot\IG(\sigma^2\,|\,\alpha_0,\beta_0),
  \label{eq:nig_joint_prior}
\end{equation}
where $\kappa_0$ sets the strength of the prior on $\mu$ relative to the data noise.

\paragraph{Derivation.}
Given data $\{x_i\}_{i=1}^n$ with sample mean $\bar{x}$ and
within-sample sum of squares $S=\sum(x_i-\bar{x})^2$, the log-posterior is
\begin{align}
  \log p(\mu,\sigma^2\,|\,\mathbf{x})
  &\propto \log p(\mathbf{x}\,|\,\mu,\sigma^2)
           + \log p(\mu\,|\,\sigma^2)
           + \log p(\sigma^2) \notag\\
  &= -\frac{n+1}{2}\log\sigma^2
     -\frac{1}{2\sigma^2}\left[
       \sum(x_i-\mu)^2 + \kappa_0(\mu-\mu_0)^2
     \right]
     - (\alpha_0+1)\log\sigma^2 - \frac{\beta_0}{\sigma^2}.
\end{align}
Completing the square in $\mu$ (treating $\sigma^2$ as fixed):
\[
  \sum(x_i-\mu)^2 + \kappa_0(\mu-\mu_0)^2
  = (\kappa_0+n)\!\left(\mu - \mu_n\right)^2 + S
  + \frac{\kappa_0 n(\bar{x}-\mu_0)^2}{\kappa_0+n},
\]
where $\mu_n = (\kappa_0\mu_0+n\bar{x})/(\kappa_0+n)$.  Substituting back:

\begin{tcolorbox}[colback=green!5, colframe=green!60!black, title=Normal-InvGamma joint update, fontupper=\small]
Posterior: $\;(\mu,\sigma^2)\,|\,\mathbf{x}\sim\mathrm{NIG}(\mu_n,\kappa_n,\alpha_n,\beta_n)$
with
\begin{align}
  \kappa_n &= \kappa_0 + n \label{eq:nig_joint_kappa}\\
  \mu_n    &= \frac{\kappa_0\mu_0 + n\bar{x}}{\kappa_n} \label{eq:nig_joint_mu}\\
  \alpha_n &= \alpha_0 + \frac{n}{2} \label{eq:nig_joint_alpha}\\
  \beta_n  &= \beta_0 + \frac{S}{2}
             + \frac{\kappa_0 n(\bar{x}-\mu_0)^2}{2\kappa_n} \label{eq:nig_joint_beta}
\end{align}
\end{tcolorbox}

\paragraph{Marginal posteriors.}
Integrating out $\sigma^2$ gives the marginal for $\mu$:
\[
  \mu\,|\,\mathbf{x}\;\sim\;
  t_{2\alpha_n}\!\left(\mu_n,\;\sqrt{\frac{\beta_n(\kappa_n+1)}{\alpha_n\kappa_n}}\right),
\]
a Student-$t$ with $2\alpha_n$ degrees of freedom, heavier-tailed than a Gaussian.
The marginal for $\sigma^2$ is simply $\IG(\alpha_n,\beta_n)$.

\paragraph{Pseudo-data interpretation.}
$\kappa_0$ imaginary observations with mean $\mu_0$ inform $\mu$;
$2\alpha_0$ imaginary observations with half-sum-of-squares $\beta_0$ inform $\sigma^2$.
The real data updates both counters simultaneously.

\begin{figure}[H]
  \centering
  \includegraphics[width=\textwidth]{fig5_joint_nig}
  \caption{Joint NIG posterior contours in $(\mu,\sigma^2)$ space.
    \textit{Left}: the prior spans a wide region.
    \textit{Centre}: after $n=8$ observations, the contours begin to
    concentrate.
    \textit{Right}: after $n=50$, the posterior is a tight ellipse
    around the true values (red star).}
  \label{fig:joint_nig}
\end{figure}

\subsection{Connection to the Darcy Gibbs step}

\begin{figure}[H]
  \centering
  \includegraphics[width=\textwidth]{fig6_darcy_gibbs}
  \caption{The InvGamma conjugate update applied to Darcy residuals.
    \textit{Left}: the prior InvGamma is flat over the relevant range;
    after conditioning on 7 sensor residuals the posterior concentrates
    near the true $\sigma_\varepsilon^2$.
    \textit{Right}: the closed-form update formula as used in the
    Gibbs step of the pCN sampler.}
  \label{fig:darcy_gibbs}
\end{figure}

Table~\ref{tab:conjugate} summarises the three conjugate pairs.

\begin{table}[H]
  \centering
  \caption{Conjugate pairs for the Gaussian likelihood.}
  \label{tab:conjugate}
  \begin{tabular}{llll}
    \toprule
    Unknown & Prior family & Posterior family & Update rule \\
    \midrule
    $\mu$ ($\sigma$ known) & Normal & Normal &
      $\tau_n^{-2} = \tau_0^{-2} + n\sigma^{-2}$ \\[2pt]
    $\sigma^2$ ($\mu$ known) & InvGamma & InvGamma &
      $\alpha_n=\alpha_0+n/2,\;\beta_n=\beta_0+\mathrm{SS}/2$ \\[2pt]
    $(\mu,\sigma^2)$ joint & Normal-InvGamma & Normal-InvGamma &
      Eqs.~\eqref{eq:nig_joint_kappa}--\eqref{eq:nig_joint_beta} \\
    \bottomrule
  \end{tabular}
\end{table}

% =====================================================================
\section{The pCN-Gibbs MCMC Sampler}
% =====================================================================

\subsection{Why not random-walk Metropolis-Hastings?}

A naive random-walk MH proposal $\xi^*=\xi+\varepsilon w$, $w\sim\N(0,I)$
suffers from the \emph{curse of dimensionality}: the optimal step size
$\varepsilon\sim M^{-1/2}$ shrinks as the number of modes $M$ grows, causing
the acceptance rate to collapse.

\subsection{The pCN proposal}

The preconditioned Crank-Nicolson (pCN) proposal is designed for targets
of the form $p(\xi\,|\,\mathbf{d})\propto p(\mathbf{d}\,|\,\xi)\,\N(\xi;0,I)$:
\begin{equation}
  \xi^* = \sqrt{1-\beta^2}\,\xi + \beta\,w, \qquad w\sim\N(0,I_M),
  \quad \beta\in(0,1).
  \label{eq:pcn}
\end{equation}

\paragraph{Prior invariance.}
If $\xi\sim\N(0,I)$ then $\xi^*\sim\N(0,I)$: the proposal is
\emph{reversible with respect to the prior}.  The Metropolis
acceptance ratio therefore only involves the likelihood ratio:
\begin{equation}
  \alpha(\xi,\xi^*)
  = \min\!\left(1,\;
    \frac{p(\mathbf{d}_{\mathrm{obs}}\,|\,\xi^*,\sigma_\varepsilon^2)}
         {p(\mathbf{d}_{\mathrm{obs}}\,|\,\xi,\sigma_\varepsilon^2)}
  \right)
  = \min\!\left(1,\;
    \exp\!\left[\frac{\norm{\mathbf{r}}^2-\norm{\mathbf{r}^*}^2}{2\sigma_\varepsilon^2}\right]\right),
  \label{eq:pcn_accept}
\end{equation}
where $\mathbf{r}=\mathbf{d}_{\mathrm{obs}}-G(\xi)$ and
$\mathbf{r}^*=\mathbf{d}_{\mathrm{obs}}-G(\xi^*)$.

\paragraph{Dimension-robustness.}
The acceptance rate depends only on the likelihood increment, not on $M$.
The same $\beta$ works regardless of the number of KL modes.

\subsection{Full sampler}

\begin{algorithm}[H]
\caption{pCN-Gibbs sampler for permeability inversion}
\KwIn{Data $\mathbf{d}_{\mathrm{obs}}$, forward model $G$, KL object $\mathcal{K}$,
      step size $\beta$, hyperparameters $(\alpha_0,\beta_0)$,
      burn-in $N_b$, samples $N_s$, thinning $\Delta$}
\KwOut{Posterior samples $\{\xi^{(t)},\sigma_\varepsilon^{2(t)}\}$}
\BlankLine
Initialise $\xi^{(0)}\leftarrow\mathbf{0}$,\;
           $\sigma_\varepsilon^{2(0)}\leftarrow 10^{-3}$,\;
           $\mathbf{d}_{\mathrm{cur}}\leftarrow G(\xi^{(0)})$\;
\For{$t = 1$ \KwTo $N_b + N_s$}{
  \textbf{Step 1 (pCN proposal for $\xi$)}\;
  Draw $w\sim\N(0,I_M)$\;
  Set $\xi^*\leftarrow\sqrt{1-\beta^2}\,\xi^{(t-1)}+\beta\,w$\;
  Evaluate $\mathbf{d}^*\leftarrow G(\xi^*)$\;
  Compute log-acceptance
    $\Delta\ell = \log p(\mathbf{d}_{\mathrm{obs}}\,|\,\xi^*,\sigma^2)
                 -\log p(\mathbf{d}_{\mathrm{obs}}\,|\,\xi^{(t-1)},\sigma^2)$\;
  Draw $u\sim\mathrm{Unif}(0,1)$\;
  \eIf{$\log u < \Delta\ell$}{
    $\xi^{(t)}\leftarrow\xi^*$;\; $\mathbf{d}_{\mathrm{cur}}\leftarrow\mathbf{d}^*$
  }{
    $\xi^{(t)}\leftarrow\xi^{(t-1)}$
  }
  \BlankLine
  \textbf{Step 2 (Gibbs draw for $\sigma_\varepsilon^2$)}\;
  Set $\mathbf{r}\leftarrow\mathbf{d}_{\mathrm{obs}}-\mathbf{d}_{\mathrm{cur}}$\;
  Draw $\sigma_\varepsilon^{2(t)}\sim\IG\!\bigl(\alpha_0+m/2,\;\beta_0+\norm{\mathbf{r}}^2/2\bigr)$
  \tcp*{Exact conjugate draw -- no MH}
  \BlankLine
  \If{$t > N_b$ \textbf{and} $(t-N_b)\bmod\Delta = 0$}{
    Store $(\xi^{(t)},\,\sigma_\varepsilon^{2(t)})$
  }
}
\end{algorithm}

\subsection{MCMC diagnostics}

\begin{figure}[H]
  \centering
  \includegraphics[width=\textwidth]{fig4_mcmc_diagnostics}
  \caption{pCN-Gibbs MCMC diagnostics.
    \textit{Left}: log-likelihood trace shows rapid mixing after the
    burn-in period.
    \textit{Centre}: noise variance $\sigma_\varepsilon^2$ trace mixes
    around the true value (red dashed).
    \textit{Right}: posterior histogram of $\sigma_\varepsilon^2$ overlaid
    with the analytical InvGamma posterior at the posterior-mean $\xi$,
    demonstrating that the Gibbs step is correctly implemented.}
  \label{fig:mcmc_diag}
\end{figure}

% =====================================================================
\section{Posterior Results}
% =====================================================================

\begin{figure}[H]
  \centering
  \includegraphics[width=\textwidth]{fig5_posterior_field}
  \caption{Posterior vs.\ true permeability field.
    \textit{Left}: log-permeability posterior mean (blue), 95\% pointwise
    credible interval (shaded), and true field (red).  The prior RMSE
    (mean field vs.\ truth) of 0.78 is reduced to 0.36 in the posterior.
    Coverage of the 95\% CI is $\approx100\%$ for this case, indicating
    slight over-coverage from the Gibbs-inferred noise variance.
    \textit{Right}: posterior predictive pressure compared to the
    noisy observations.  The uncertainty in $k(x)$ propagates to
    uncertainty in $p(x)$.}
  \label{fig:posterior}
\end{figure}

\begin{figure}[H]
  \centering
  \includegraphics[width=0.72\textwidth]{fig6_posterior_predictive}
  \caption{Posterior predictive check in data space.
    Blue squares show the posterior predictive mean $\pm2\sigma$;
    red triangles are the noisy observations.
    The posterior predictive brackets the observations, confirming
    that the inferred field is consistent with the experiment.}
  \label{fig:pred_check}
\end{figure}

\subsection{Credible intervals versus confidence intervals}

A 95\% \emph{credible interval} $[a,b]$ satisfies
$P(k(x)\in[a,b]\,|\,\mathbf{d}_{\mathrm{obs}})=0.95$---a direct
probability statement about the unknown quantity given the data.
This is distinct from a frequentist 95\% confidence interval, which
would state that 95\% of \emph{repeated experiments} would produce
intervals containing the true value.  In the UQ context, the Bayesian
credible interval is the natural quantity: we have \emph{one}
experiment and wish to quantify uncertainty given that experiment.

% =====================================================================
\section{Conclusions}
% =====================================================================

We have traced the full chain from textbook probability distributions to
a working posterior sampler for 1D single-phase Darcy permeability
inversion:

\begin{enumerate}
  \item The \textbf{Normal distribution} enters as the Gaussian likelihood
    (measurement error model), the GRF prior on $\log k(x)$, and the
    KL coefficients.

  \item The \textbf{Inverse-Gamma} enters as the conjugate hyperprior
    on $\sigma_\varepsilon^2$, enabling an exact Gibbs update that
    replaces a Metropolis step.

  \item The \textbf{Gamma} is the dual form (prior on precision),
    equally valid and leading to identical posterior calculations.

  \item The \textbf{KL expansion} compresses the infinite-dimensional
    field inference into a finite-dimensional problem in $\R^M$, with
    the prior $\xi\sim\N(0,I)$ enabling the dimension-robust pCN proposal.

  \item The \textbf{pCN-Gibbs} sampler alternates between a
    prior-reversible Metropolis step for $\xi$ and an exact conjugate
    draw for $\sigma_\varepsilon^2$, producing samples from the full
    joint posterior.
\end{enumerate}

The posterior log-$k$ RMSE (0.36) is less than half the prior RMSE (0.78),
and the 95\% credible intervals comfortably cover the true field---concrete
evidence that the 7 sensor readings carry meaningful information about the
permeability structure.

% =====================================================================
\appendix
\section{Complete Derivation of the NIG Joint Update}
% =====================================================================

Starting from the log-posterior:
\[
  \log p(\mu,\sigma^2\,|\,\mathbf{x})
  \propto
  -\frac{n+1}{2}\log\sigma^2
  -\frac{1}{2\sigma^2}\!\left[
    \sum(x_i-\mu)^2 + \kappa_0(\mu-\mu_0)^2 + 2\beta_0
  \right].
\]
Define $\kappa_n=\kappa_0+n$ and $\mu_n=(\kappa_0\mu_0+n\bar{x})/\kappa_n$.
Expanding the quadratic in $\mu$:
\begin{align*}
  \sum(x_i-\mu)^2 + \kappa_0(\mu-\mu_0)^2
  &= \sum x_i^2 - 2n\bar{x}\mu + n\mu^2
     + \kappa_0\mu^2 - 2\kappa_0\mu_0\mu + \kappa_0\mu_0^2 \\
  &= \kappa_n\mu^2 - 2\kappa_n\mu_n\mu + (n\bar{x}^2+\kappa_0\mu_0^2-\kappa_n\mu_n^2) + S \\
  &= \kappa_n(\mu-\mu_n)^2 + S + \frac{\kappa_0 n(\bar{x}-\mu_0)^2}{\kappa_n},
\end{align*}
where $S=\sum(x_i-\bar{x})^2$ and the last equality uses
$n\bar{x}^2+\kappa_0\mu_0^2-\kappa_n\mu_n^2=\kappa_0 n(\bar{x}-\mu_0)^2/\kappa_n$
(verified by direct expansion).

Substituting:
\[
  \log p(\mu,\sigma^2\,|\,\mathbf{x})
  \propto
  -\frac{\kappa_n(\mu-\mu_n)^2}{2\sigma^2}
  -\left(\alpha_0+\frac{n}{2}+1\right)\!\log\sigma^2
  -\frac{\beta_0+S/2+\kappa_0 n(\bar{x}-\mu_0)^2/(2\kappa_n)}{\sigma^2}.
\]
This is the log of $\N(\mu\,|\,\mu_n,\sigma^2/\kappa_n)\cdot\IG(\sigma^2\,|\,\alpha_n,\beta_n)$
with $\alpha_n,\beta_n$ as given in
Eqs.~\eqref{eq:nig_joint_alpha}--\eqref{eq:nig_joint_beta}. \hfill$\square$

% =====================================================================
\section{The Gamma Distribution as Conjugate Prior for Precision:
         Full Algebraic Derivation}
\label{app:gamma_precision}
% =====================================================================

This appendix derives in full why $\Ga(\alpha_0,\beta_0)$ placed on the
Gaussian \emph{precision} $\tau=1/\sigma^2$ produces a Gamma posterior.
Every algebraic step is explicit.

\subsection*{B.1\quad Setup}

\paragraph{Model.}
\[
  x_i\,|\,\tau
    \overset{\mathrm{i.i.d.}}{\sim}
    \N\!\left(\mu,\,\frac{1}{\tau}\right),
  \quad i=1,\ldots,n,
  \qquad
  \mu \text{ known}, \quad \tau>0 \text{ unknown}.
\]
\paragraph{Prior.}
\[
  \tau \sim \Ga(\alpha_0,\beta_0)
  \qquad\Longrightarrow\qquad
  p(\tau;\alpha_0,\beta_0)
    = \frac{\beta_0^{\alpha_0}}{\Gamma(\alpha_0)}\,
      \tau^{\alpha_0-1}\,e^{-\beta_0\tau},
  \quad \tau>0.
\]
Here $\beta_0$ is the \emph{rate} (inverse scale), so
$\E[\tau]=\alpha_0/\beta_0$ and $\mathrm{Mode}[\tau]=(\alpha_0-1)/\beta_0$
for $\alpha_0>1$.

\subsection*{B.2\quad Gaussian density rewritten in precision form}

Starting from $\N(x;\mu,\sigma^2)$ and substituting $\tau=1/\sigma^2$:
\begin{equation}
  \N\!\left(x;\mu,\tfrac{1}{\tau}\right)
  = \frac{1}{\sqrt{2\pi/\tau}}\exp\!\left(-\frac{(x-\mu)^2}{2/\tau}\right)
  = \sqrt{\frac{\tau}{2\pi}}\;\exp\!\left(-\frac{\tau(x-\mu)^2}{2}\right).
  \label{eq:gauss_prec}
\end{equation}
The variance $\sigma^2$ in the denominator has become the precision $\tau$
in the numerator of the exponent.  This is the key structural fact: the
likelihood is \textbf{linear in $\tau$ in the exponent}.

\subsection*{B.3\quad The likelihood}

\begin{align}
  p(\mathbf{x}\,|\,\tau)
  &= \prod_{i=1}^n \sqrt{\frac{\tau}{2\pi}}\,
     \exp\!\left(-\frac{\tau(x_i-\mu)^2}{2}\right) \notag\\
  &= (2\pi)^{-n/2}\;\tau^{n/2}\;
     \exp\!\left(-\frac{\tau}{2}\sum_{i=1}^n(x_i-\mu)^2\right) \notag\\
  &= (2\pi)^{-n/2}\;\tau^{n/2}\;\exp\!\left(-\frac{\mathrm{SS}}{2}\,\tau\right),
  \label{eq:lik_prec}
\end{align}
where $\mathrm{SS}=\sum_{i=1}^n(x_i-\mu)^2$ is the sum of squared
deviations from the known mean.

The $\tau$-dependent part of \eqref{eq:lik_prec} is
\[
  \tau^{n/2}\,e^{-(\mathrm{SS}/2)\,\tau}.
\]
Compare this with the Gamma kernel $\tau^{\alpha_0-1}\,e^{-\beta_0\tau}$:
both are a power of $\tau$ times $\exp(-\mathrm{const}\cdot\tau)$.

\subsection*{B.4\quad Posterior derivation (kernel recognition)}

Applying Bayes' theorem and dropping all $\tau$-free constants:
\begin{align}
  p(\tau\,|\,\mathbf{x})
  &\propto p(\mathbf{x}\,|\,\tau)\cdot p(\tau) \notag\\[4pt]
  &\propto \Bigl[\tau^{n/2}\,e^{-(\mathrm{SS}/2)\tau}\Bigr]
           \;\cdot\;
           \Bigl[\tau^{\alpha_0-1}\,e^{-\beta_0\tau}\Bigr] \notag\\[4pt]
  &= \tau^{\,n/2\;+\;\alpha_0\;-\;1}
     \cdot\exp\!\bigl(-\underbrace{(\beta_0+\mathrm{SS}/2)}_{\displaystyle\beta_n}\,\tau\bigr)
     \notag\\[4pt]
  &= \tau^{\,\underbrace{(\alpha_0+n/2)}_{\displaystyle\alpha_n}-1}
     \cdot e^{-\beta_n\tau}.
  \label{eq:gamma_post_kernel}
\end{align}

Expression \eqref{eq:gamma_post_kernel} is the \emph{unnormalised kernel}
of $\Ga(\alpha_n,\beta_n)$.  Since the Gamma normalising constant
$\beta_n^{\alpha_n}/\Gamma(\alpha_n)$ is uniquely determined by
$(\alpha_n,\beta_n)$, we conclude immediately:

\begin{tcolorbox}[colback=blue!5, colframe=blue!50!black,
  title=Gamma-precision conjugate update, fontupper=\small]
\[
  \tau\,|\,\mathbf{x}_1,\ldots,\mathbf{x}_n
  \;\sim\;
  \Ga(\alpha_n,\beta_n)
\]
\vspace{-6pt}
\begin{align}
  \alpha_n &= \alpha_0 + \frac{n}{2}
    \tag{shape: $\tfrac{1}{2}$ added per observation} \\[4pt]
  \beta_n  &= \beta_0 + \frac{\mathrm{SS}}{2}
    = \beta_0 + \frac{1}{2}\sum_{i=1}^n(x_i-\mu)^2
    \tag{rate: squared deviations accumulate}
\end{align}
\end{tcolorbox}

\subsection*{B.5\quad Why the exponents collapse: exponential family structure}

Both distributions belong to the \emph{exponential family} in the natural
parameter $\eta=-\tau$ (or equivalently, with sufficient statistic $T(\tau)=\tau$):

\medskip
\begin{center}
\begin{tabular}{llll}
\toprule
Distribution & Kernel & Natural param $\eta$ & Suff.\ stat $T$ \\
\midrule
$\Ga(\alpha_0,\beta_0)$ prior   & $\tau^{\alpha_0-1}e^{-\beta_0\tau}$ & $-\beta_0$ & $\tau$ \\
Gaussian likelihood              & $\tau^{n/2}e^{-(\mathrm{SS}/2)\tau}$ & $-\mathrm{SS}/2$ & $\tau$ \\
\midrule
Product (posterior kernel)       & $\tau^{(\alpha_0+n/2)-1}e^{-(\beta_0+\mathrm{SS}/2)\tau}$ & $-(\beta_0+\tfrac{\mathrm{SS}}{2})$ & $\tau$ \\
\bottomrule
\end{tabular}
\end{center}
\medskip

Multiplying two exponentials with the same sufficient statistic
simply \emph{adds} the natural parameters:
\[
  e^{-\beta_0\tau}\cdot e^{-(\mathrm{SS}/2)\tau}
  = e^{-(\beta_0+\mathrm{SS}/2)\tau}.
\]
This is the algebraic essence of conjugacy: the prior and likelihood share
the same exponential structure in the parameter $\tau$, so their product
stays in the same exponential family.

\subsection*{B.6\quad Normalising constant and marginal likelihood}

The full posterior normalisation gives the \emph{marginal likelihood}
$p(\mathbf{x})$ for free:
\begin{align}
  p(\mathbf{x})
  &= \int_0^\infty p(\mathbf{x}\,|\,\tau)\,p(\tau)\,\mathrm{d}\tau \notag\\
  &= (2\pi)^{-n/2}\,\frac{\beta_0^{\alpha_0}}{\Gamma(\alpha_0)}
     \int_0^\infty \tau^{\alpha_n-1}\,e^{-\beta_n\tau}\,\mathrm{d}\tau \notag\\
  &= (2\pi)^{-n/2}\,\frac{\beta_0^{\alpha_0}}{\Gamma(\alpha_0)}
     \cdot\frac{\Gamma(\alpha_n)}{\beta_n^{\alpha_n}},
  \label{eq:marglik}
\end{align}
using $\int_0^\infty \tau^{\alpha_n-1}e^{-\beta_n\tau}\,\mathrm{d}\tau
= \Gamma(\alpha_n)/\beta_n^{\alpha_n}$.
Equation~\eqref{eq:marglik} is a product of Gamma and Student-t terms
and is the quantity used in Bayesian model comparison.

\subsection*{B.7\quad Sequential consistency}

Suppose we process the $n$ observations one at a time.  After datum $x_1$:
\[
  \alpha_1 = \alpha_0 + \tfrac{1}{2}, \qquad
  \beta_1  = \beta_0  + \tfrac{(x_1-\mu)^2}{2}.
\]
After datum $x_2$ (starting from $\Ga(\alpha_1,\beta_1)$):
\[
  \alpha_2 = \alpha_1 + \tfrac{1}{2} = \alpha_0 + 1, \qquad
  \beta_2  = \beta_1  + \tfrac{(x_2-\mu)^2}{2}
           = \beta_0  + \tfrac{(x_1-\mu)^2+(x_2-\mu)^2}{2}.
\]
After all $n$ data:
\[
  \alpha_n = \alpha_0 + \frac{n}{2}, \qquad
  \beta_n  = \beta_0 + \frac{\sum_{i=1}^n(x_i-\mu)^2}{2}
           = \beta_0 + \frac{\mathrm{SS}}{2}.
\]
This is identical to the batch update---the order of data presentation is irrelevant.
This is sequentially consistent Bayesian updating.

\subsection*{B.8\quad Pseudo-data interpretation}

The prior $\Ga(\alpha_0,\beta_0)$ can be viewed as arising from
$n_0=2\alpha_0$ \emph{imaginary} (pseudo-)~observations whose
sum of squared deviations is $\mathrm{SS}_0=2\beta_0$.  Then the
posterior after $n$ real observations is
\[
  \tau\,|\,\mathbf{x}
  \;\sim\;
  \Ga\!\left(\alpha_0+\frac{n}{2},\;
              \beta_0+\frac{\mathrm{SS}}{2}\right)
  \equiv
  \Ga\!\left(\frac{n_0+n}{2},\;
              \frac{\mathrm{SS}_0+\mathrm{SS}}{2}\right),
\]
exactly as if we had $n_0+n$ total observations with combined sum of
squares $\mathrm{SS}_0+\mathrm{SS}$.  The prior encodes the same type
of information as data; real data simply extends the pseudo-dataset.

\subsection*{B.9\quad The dual InvGamma via change of variables}

Since $\tau>0$, define $\sigma^2=1/\tau$.  The Jacobian is
$\mathrm{d}\tau/\mathrm{d}(\sigma^2) = -1/(\sigma^2)^2 = -\sigma^{-4}$, so
\begin{align}
  p(\sigma^2)
  &= p\!\left(\tau=\frac{1}{\sigma^2}\right)\cdot\left|\frac{\mathrm{d}\tau}{\mathrm{d}(\sigma^2)}\right|
   = \frac{\beta^{\alpha}}{\Gamma(\alpha)}
     \left(\frac{1}{\sigma^2}\right)^{\!\alpha-1}
     e^{-\beta/\sigma^2}\cdot\frac{1}{\sigma^4} \notag\\
  &= \frac{\beta^{\alpha}}{\Gamma(\alpha)}
     \left(\sigma^2\right)^{-\alpha-1}
     e^{-\beta/\sigma^2}.
  \label{eq:invgamma_derivation}
\end{align}
This is exactly $\IG(\alpha,\beta)$.  Therefore:
\[
  \tau\sim\Ga(\alpha,\beta)
  \;\;\Longleftrightarrow\;\;
  \sigma^2=\frac{1}{\tau}\sim\IG(\alpha,\beta).
\]
The update formulas are identical in both parametrisations
($\alpha_n=\alpha_0+n/2$, $\beta_n=\beta_0+\mathrm{SS}/2$);
only the axis label changes from $\tau$ to $\sigma^2=1/\tau$.

\medskip
\begin{center}
\begin{tabular}{lll}
\toprule
Parametrisation & Prior & Posterior \\
\midrule
Precision $\tau=1/\sigma^2$ & $\Ga(\alpha_0,\beta_0)$ & $\Ga(\alpha_0+n/2,\;\beta_0+\mathrm{SS}/2)$ \\
Variance $\sigma^2=1/\tau$  & $\IG(\alpha_0,\beta_0)$ & $\IG(\alpha_0+n/2,\;\beta_0+\mathrm{SS}/2)$ \\
\bottomrule
\end{tabular}
\end{center}

\subsection*{B.10\quad Summary of the conjugacy argument}

The proof reduces to one observation: the Gaussian likelihood
$p(\mathbf{x}|\tau)\propto\tau^{n/2}e^{-(\mathrm{SS}/2)\tau}$ and
the Gamma prior $p(\tau)\propto\tau^{\alpha_0-1}e^{-\beta_0\tau}$
are both of the form $\tau^{c}\cdot e^{-d\tau}$.  Multiplying two such
expressions yields another expression of the same form with
updated exponent $c$ and rate $d$:
\[
  \underbrace{\tau^{\alpha_0-1}\,e^{-\beta_0\tau}}_{\text{prior}}
  \;\times\;
  \underbrace{\tau^{n/2}\,e^{-(\mathrm{SS}/2)\tau}}_{\text{likelihood}}
  \;=\;
  \tau^{(\alpha_0+n/2)-1}\,e^{-(\beta_0+\mathrm{SS}/2)\tau}
  \;\;\propto\;\;
  \underbrace{\Ga\!\left(\alpha_0+\tfrac{n}{2},\;\beta_0+\tfrac{\mathrm{SS}}{2}\right)}_{\text{posterior}}.
\]
The family is preserved because the product of two power-times-exponential
functions is again a power-times-exponential function---and that is
precisely the Gamma family.

\end{document}
